
\begin{frame}{Appendix --- Pre-Trend Placebo (consolidation ladder, numeric)}
\hypertarget{app:pretrend}{}
\small
The numeric placebo behind the balance coefplot: the 2016$\to$2020 pre-period change in each consolidation outcome, regressed on the instrumented 2020$\to$2024 treatment. A future shock cannot cause a past change, so a clean design returns coefficients near zero. Conditioned on pre-determined (2010 Census) covariates only; state FE; state-clustered.
\smallskip
\begin{center}
\scriptsize
\resizebox{\linewidth}{!}{% Auto-generated by code/03_estimation/02_iv_main.R
% Do not edit manually -- rerun the generating script to update

\begin{tabular}{l*{4}{c}}
\toprule\toprule
Dep.\ var.: & \textbf{Margin (W$-$RU)} & \textbf{Eff.\ N candidates} & \textbf{Vote HHI} & \textbf{Top-2 vote share} \\
\midrule
\rowcolor{mylight}
\textbf{Placebo (judicialization)} & 0.003 & -0.035 & 0.013 & 0.011 \\
\rowcolor{mylight}
 & \textcolor{mygray}{(0.038)} & \textcolor{mygray}{(0.070)} & \textcolor{mygray}{(0.018)} & \textcolor{mygray}{(0.012)} \\
\midrule
$N$ & 5,560 & 5,560 & 5,560 & 5,560 \\
Pre-trend mean & 0.017 & 0.087 & -0.011 & -0.022 \\
\bottomrule\bottomrule
\end{tabular}
}
\end{center}
\smallskip
\footnotesize Every rung is clean: margin ($\PtMarginCoef$, $p=\PtMarginP$), effective-N ($\PtEffNCoef$, $p=\PtEffNP$), Herfindahl ($\PtHHICoef$, $p=\PtHHIP$), top-two share ($\PtTopTwoCoef$, $p=\PtTopTwoP$). This differs from an earlier draft because that version conditioned the balance test on \texttt{margin\_2016}, the ladder's own 2016 baseline, which induced a regression-to-the-mean (Lord's paradox) artifact that read as a pre-trend. The baseline is inadmissible when the outcome is itself a 2016-based change, and dropping it clears every rung. \hfill\hyperlink{src:pretrend}{\beamerreturnbutton{Back to balance test}}
\end{frame}
